% paper.tex -- draft, A&A LaTeX class (aa.cls v7.0, EDP Sciences)
% Purely toroidal branch: theory, self-consistent-field method, and results.
% TODO: author list, affiliations, figures.

\documentclass{aa}

\usepackage{graphicx}
\usepackage{amsmath}
\usepackage{amssymb}
\usepackage{txfonts}
\usepackage{natbib}

\begin{document}

   \title{Super-Chandrasekhar white dwarfs supported by a purely
          toroidal magnetic field}

   \subtitle{Draft}

   \author{TBD\inst{1}
          }

   \institute{Affiliation TBD\\
              \email{tbd@example.org}
             }

   \date{Draft \today}

   \abstract
   % context
   {A white dwarf mass appreciably above the Chandrasekhar limit is
    usually attributed to rotation. A strong internal magnetic field is
    the other candidate, and the purely toroidal configuration is the
    one for which the equilibrium problem remains tractable without a
    Grad-Shafranov solve.}
   % aims
   {We quantify how much equilibrium mass a purely toroidal field adds
    to a cold, non-rotating white dwarf, and how that mass gain depends
    on the central density, within the density range in which the cold
    degenerate equation of state can still be claimed to apply.}
   % methods
   {The field is prescribed as $B_\varphi = K\rho^m\varpi^{2m-1}$,
    the form for which the Lorentz force derives from a potential and
    the barotropic structure equations close. The resulting
    magneto-hydrostatic equilibria are computed with a self-consistent
    field iteration in which the density inversion becomes an implicit,
    per-point root find, and are certified with the virial error and
    with an explicit test that the configuration fits inside its
    computational domain.}
   % results
   {The dimensionless field strength $E_{\mathrm{mag}}/|W|$ is predicted
    analytically to be independent of central density and radius, and is
    measured to vary by $0.9\%$ across two and a half decades in
    $\rho_c$. A field with $E_{\mathrm{tor}}/|W| \simeq 0.18$ raises the
    equilibrium mass by $\simeq 47\%$, essentially independent of
    central density, giving $M > 2\,M_\odot$ inside the density range
    allowed by the inverse-beta-decay threshold, with a maximum interior
    field of $2$--$3\times10^{14}$\,G. Adding rigid rotation on top of
    such a field shortens, rather than extends, the range of angular
    velocity over which the sequence remains virially certified.}
   % conclusions
   {A purely toroidal field is sufficient, on its own, to reach
    $2\,M_\odot$ at physically admissible central densities. It is,
    however, the configuration most exposed to the Tayler $m=1$
    instability, which the axisymmetric formulation used here cannot
    test, and it produces no exterior field and hence no dipole for
    magnetic braking to act on.}

   \keywords{white dwarfs -- stars: magnetic field -- magnetohydrodynamics (MHD) -- methods: numerical}

   \maketitle

%-------------------------------------------------------------------

\section{Introduction}

TBD.

%-------------------------------------------------------------------

\section{Equation of state}
\label{sec:eos}

The background stellar structure throughout this work is that of a
fully degenerate, zero-temperature electron gas
\citep[ztwd;][]{chandrasekhar1935}. Pressure and density are both
parametrized by the normalized Fermi momentum $x \equiv p_F/(m_e c)$,

\begin{align}
P(x) &= A \left[ x(2x^2-3)\sqrt{x^2+1} + 3\sinh^{-1}x \right], \\
\rho(x) &= B\,x^3,
\end{align}

\noindent where $A$ is a fixed combination of fundamental constants
($m_e$, $c$, $\hbar$) and $B$ depends on the composition only through
the mean molecular weight per electron $\mu_e = 1/Y_e$, with $Y_e$ the
number of electrons per baryon. This is a purely mechanical equation of
state: at fixed composition, pressure depends on density alone, with no
temperature dependence. Consequently, once the configuration is
restricted to be non-rotating and field-free, the entire stellar
structure (mass, radius, density profile) is determined by exactly two
physical parameters: the central density $\rho_c$ and $\mu_e$.

The specific enthalpy has a closed form,
\begin{equation}
H(x) = \frac{8A}{B}\left[\sqrt{1+x^2} - 1\right],
\label{eq:enthalpy}
\end{equation}
which is what makes the self-consistent field (SCF) iteration of
Sect.~\ref{sec:scf} tractable: in the field-free case the density is
recovered directly from the local enthalpy via the inverse relation
$x(H) = \sqrt{(1+HB/8A)^2 - 1}$, with $H \le 0$ mapped to $\rho = 0$
(the stellar surface). Section~\ref{sec:implicit} describes how this
step changes once a toroidal field is present.

Near the surface ($x \ll 1$) the equation of state reduces to
$\rho \propto H^{3/2}$, an effective polytropic index $n=3/2$; in the
relativistic core ($x \gg 1$) it approaches $\rho \propto H^3$,
$n=3$ -- the marginal polytropic index at which the equilibrium mass
becomes independent of central density, the origin of the
Chandrasekhar mass limit.

At sufficiently high density, inverse beta decay (electron capture)
changes the composition locally and the constant-$\mu_e$ cold equation
of state above ceases to describe a real white dwarf. We adopt the
tabulated threshold densities of \citet{boshkayev2013} for this
validity limit. This gate is not a formality here: it is what
distinguishes a mass gain that can be claimed for a white dwarf from
one obtained at a central density where the assumed microphysics no
longer holds, and every configuration quoted as a result in
Sect.~\ref{sec:results} lies below the threshold for the assumed
composition.

%-------------------------------------------------------------------

\section{A purely toroidal field}
\label{sec:toroidal}

\subsection{The field law}

For a barotropic star the magnetic force must derive from a potential
if hydrostatic balance is to close in the form of a Bernoulli integral.
For an axisymmetric, purely azimuthal field
$\mathbf{B} = B_\varphi(\varpi,z)\,\hat{\mathbf{e}}_\varphi$, where
$\varpi$ is the cylindrical radius, this restricts $B_\varphi$ to the
one-parameter family
\begin{equation}
B_\varphi = K \rho^m \varpi^{2m-1}, \qquad m \ge 1,
\label{eq:fieldlaw}
\end{equation}
with $K$ setting the field amplitude and $m$ its concentration towards
the dense inner regions. The corresponding Lorentz force, obtained from
the divergence of the Maxwell stress tensor, is
$\mathbf{f}_L/\rho = -\nabla M_{\mathrm{tor}}$ with
\begin{equation}
M_{\mathrm{tor}}(s) = \frac{m K^2}{4\pi(2m-1)}\, s^{\,2m-1},
\qquad s \equiv \rho\varpi^2 ,
\label{eq:mtor}
\end{equation}
which we verified symbolically for $m = 1, 3/2, 2, 3$ (the residual of
the force balance vanishes identically). Hydrostatic equilibrium then
reduces to the algebraic Bernoulli integral
\begin{equation}
H + \Phi + M_{\mathrm{tor}}(\rho\varpi^2) = C ,
\label{eq:bernoulli}
\end{equation}
with $\Phi$ the gravitational potential, solved together with the
Poisson equation for $\Phi$ and with the equation of state of
Sect.~\ref{sec:eos}.

Two features of Eq.~\ref{eq:fieldlaw} are worth making explicit,
because they set both the appeal and the limits of this branch. First,
$B_\varphi$ is an \emph{algebraic} function of the local density: no
flux function, no Grad-Shafranov operator, and no Green's function are
involved in obtaining the field, in contrast with the poloidal case.
Second, and for the same reason, the field vanishes wherever the
density does. A purely toroidal field is strictly confined to the
stellar interior; there is no exterior field and no dipole moment, a
point we return to in Sect.~\ref{sec:limitations}.

\subsection{The implicit density inversion}
\label{sec:implicit}

Because $M_{\mathrm{tor}}$ depends on $\rho$ -- the very quantity being
solved for -- the density update of the SCF iteration is no longer a
direct inversion of the equation of state. At each grid point one must
instead solve
\begin{equation}
H(\rho) + M_{\mathrm{tor}}(\rho\varpi^2) = C - \Phi
\end{equation}
for $\rho$. The left-hand side is monotonically increasing in $\rho$
(both $H$ and $M_{\mathrm{tor}}$ are), so the root is unique and safely
bracketed, and we obtain it with a standard bracketing root finder. This
is a change in the structure of the algorithm, not in its physics, and
it is the only place where the toroidal branch is more expensive than
the field-free one.

\subsection{Virial diagnostic}
\label{sec:virial}

Every configuration is certified with the scalar virial identity, which
for a static magnetized star reads
\begin{equation}
W + 3\Pi + \mathcal{M} = 0, \qquad
\mathrm{VE} \equiv \frac{|W + 3\Pi + \mathcal{M}|}{|W|} ,
\label{eq:ve}
\end{equation}
with $W$ the gravitational energy, $\Pi = \int P\,dV$, and
$\mathcal{M} = \int B_\varphi^2/8\pi\,dV$ the magnetic energy. For the
field law of Eq.~\ref{eq:fieldlaw} the magnetic contribution satisfies
\begin{equation}
\int \rho\, \nabla M_{\mathrm{tor}}\!\cdot\mathbf{r}\; dV
 \;=\; -\int \frac{B_\varphi^2}{8\pi}\, dV ,
\end{equation}
the sign following from $M_{\mathrm{tor}}$ entering
Eq.~\ref{eq:bernoulli} with a positive sign, combined with the
convention-free Maxwell-stress identity
$\int \mathbf{r}\cdot\mathbf{f}_L\,dV = +\int B^2/8\pi\,dV$, valid for
any purely toroidal field. We confirmed this numerically: the residual
of Eq.~\ref{eq:ve} falls monotonically, $2.27\% \rightarrow 1.00\%
\rightarrow 0.36\% \rightarrow 0.14\%$, as the radial mesh is refined
from $n_r = 65$ to $n_r = 257$. We take $\mathrm{VE} < 10^{-3}$ as the
acceptance threshold throughout, and quote it for every configuration
reported below.

%-------------------------------------------------------------------

\section{Numerical method}
\label{sec:scf}

Equilibria are obtained with the self-consistent field scheme of
\citet{hachisu1986}, in which the density field, the gravitational
potential and the Bernoulli constant are updated in turn until the
density field stops changing. The gravitational potential is obtained
from a Legendre (multipole) expansion of the Poisson Green's function,
truncated at $l_{\max}$; the magnetic term enters only through
Eq.~\ref{eq:mtor} and the implicit inversion of
Sect.~\ref{sec:implicit}. Three numerical issues proved to matter
quantitatively for this branch, and we record them because each one, on
its own, is capable of producing a wrong physical conclusion.

\paragraph{Domain size.} A toroidal field makes the star \emph{prolate}
($R_{\mathrm{pol}} > R_{\mathrm{eq}}$), the opposite of the oblate
deformation produced by rotation, and the deformation grows with $K$.
A computational domain sized from a field-free radius estimate is
therefore progressively too small as $K$ increases, and the "surface"
the solver finds is the edge of the domain rather than a physical
radius. We monitor this directly through
$f_{\mathrm{pol}} \equiv R_{\mathrm{pol}}/r_{\max}$ and accept only
configurations with $f_{\mathrm{pol}} \lesssim 0.2$; a value of $1.0$
means the star does not fit in its box at all. This is not a
conservative margin added for safety -- an earlier pass of this work
concluded that $2\,M_\odot$ was \emph{not} reachable, a conclusion
produced entirely by domain overflow at every point of the sequence
with the strongest field, and reversed once the domain was enlarged.
Domain size and mesh resolution must be varied one at a time: enlarging
the domain at fixed $\Delta r$ drops VE from $1.43\times10^{-2}$ to
$2.49\times10^{-4}$ and then leaves $M$ and VE unchanged to four
significant figures under a further enlargement, whereas refining the
mesh inside a domain that is too small does not improve VE at all.

\paragraph{Continuation in $K$.} Direct solves at $K \ge 5\times10^{-3}$
do not converge at any domain size we tried. These configurations are
reachable only by continuation: stepping $K$ upward in small increments
from the field-free solution and warm-starting each step from the
previous one.

\paragraph{The multipole expansion.} Forming the radial Green's
function as separate powers $r'^{\,l+2}$ and $r^{-(l+1)}$ before
dividing overflows double precision outright at $l = 48$ for stellar
radii, and silently loses precision even at $l_{\max} = 16$, where the
two factors are separated by some $270$ orders of magnitude before a
division intended to cancel most of that scale. Rewriting the radial
integrals as a recursion that only ever forms the ratios
$(r_{i-1}/r_i)^{l+1} \le 1$ and $(r_i/r_{i+1})^{l} \le 1$ -- the same
integral, with no absolute power of $r$ ever formed -- removes the
problem. With this in place, the truncation itself is not a limitation
for the configurations considered here: stepping
$l_{\max} = 16 \rightarrow 32 \rightarrow 48$ moves VE by less than
$10\%$.

%-------------------------------------------------------------------

\section{Results}
\label{sec:results}

\subsection{An invariance predicted before it was measured}
\label{sec:scaling}

For $m=1$ the field law is $B_\varphi = K\rho\varpi$, and the
dimensionless field strength can be estimated without solving anything.
With $E_{\mathrm{mag}} \sim K^2\rho^2R^5/8\pi$ and
$|W| \sim GM^2/R \sim G\rho^2R^5$, the density and radius cancel:
\begin{equation}
\frac{E_{\mathrm{mag}}}{|W|} \;\sim\; \frac{K^2}{8\pi G} ,
\label{eq:scaling}
\end{equation}
independent of both $\rho_c$ and $R$, up to a shape factor of order
unity that drifts slowly along the sequence as the effective polytropic
index slides from $3/2$ towards $3$. Equation~\ref{eq:scaling} predicts
that the fractional mass gain produced by a given $K$ should be
essentially the same at any central density -- and, for
$K = 3\times10^{-3}$, that it should be of order $50\%$.

Table~\ref{tab:masses} shows the measured values at three central
densities spanning two and a half decades. The prediction was made
before these runs and is confirmed by them: $E_{\mathrm{tor}}/|W|$
varies by $0.9\%$ over the full range, and the fractional mass gain by
less than three percentage points.

\begin{table}
\caption{Equilibrium mass with and without a purely toroidal field
($m=1$, $K = 3\times10^{-3}$, no rotation), at three central densities.
All points are certified at $f_{\mathrm{pol}} \simeq 0.09$--$0.16$ and
reached by continuation from $K=0$.}
\label{tab:masses}
\centering
\begin{tabular}{l c c c c c}
\hline\hline
$\rho_c$ & $M(K=0)$ & $M(K)$ & gain & $E_{\mathrm{tor}}/|W|$ & VE \\
(g\,cm$^{-3}$) & ($M_\odot$) & ($M_\odot$) & & & \\
\hline
$1.0\times10^{10}$ & 1.4108 & 2.0722 & $+46.9\%$ & 0.1809 & $4.0\times10^{-4}$ \\
$1.5\times10^{10}$ & 1.4159 & 2.0874 & $+47.4\%$ & 0.1805 & $3.9\times10^{-4}$ \\
$1.0\times10^{12}$ & 1.4323 & 2.1426 & $+49.6\%$ & 0.1790 & $2.7\times10^{-4}$ \\
\hline
\end{tabular}
\tablefoot{The $\rho_c = 10^{12}$\,g\,cm$^{-3}$ row is given for
reference only: it lies above the inverse-beta-decay threshold
(Sect.~\ref{sec:eos}) and is not a physically admissible white dwarf.}
\end{table}

\subsection{The mass gain}

The result is best stated in terms of $E_{\mathrm{tor}}/|W|$ rather
than $K$, since $K$ carries the units and the arbitrariness of the
ansatz of Eq.~\ref{eq:fieldlaw} and means nothing outside it, whereas
the energy ratio transfers to any other field configuration: a toroidal
field with $E_{\mathrm{tor}}/|W| \simeq 0.18$ raises the equilibrium
mass by $\simeq 47\%$, essentially independent of central density,
which puts $M > 2\,M_\odot$ inside the density range the
inverse-beta-decay threshold allows. The two admissible rows of
Table~\ref{tab:masses} both exceed $2\,M_\odot$ while remaining below
the neutronization threshold for $\mu_e = 2$; the mass gain does not
have to be bought by pushing $\rho_c$ into a regime the equation of
state cannot describe.

In physical units, the corresponding interior fields at those two
points are
\begin{equation*}
B_{\varphi,\max} = 1.95\times10^{14}\ \mathrm{G} \quad\text{and}\quad
2.56\times10^{14}\ \mathrm{G},
\end{equation*}
with volume-weighted means of $7.9\times10^{12}$ and
$9.9\times10^{12}$\,G respectively. These are the field strengths that
come out of the calculation, with no tuning towards any target value.

\subsection{Rotation does not extend the sequence}
\label{sec:rotation}

Since a toroidal field elongates the star and rotation flattens it, a
natural expectation is that combining the two should push the
mass-shedding limit to higher angular velocity than rotation alone. We
tested this directly, running the same continuation in the central
angular velocity $\Omega_c$ twice on an identical mesh and at identical
$\rho_c$, once field-free and once with a moderate toroidal field
($K = 3\times10^{-3}$, $E_{\mathrm{tor}}/|W| \simeq 0.18$ at
$\Omega_c = 0$).

The expectation fails. With the field present, the equatorial mass-loss
ratio and the virial error degrade \emph{together} and smoothly between
$\Omega_c = 15$ and $20$\,rad\,s$^{-1}$: the mass-loss ratio crosses
unity, and VE crosses the $10^{-3}$ acceptance threshold at essentially
the same $\Omega_c$ and then climbs monotonically to $0.16$ by
$\Omega_c = 32$\,rad\,s$^{-1}$. The practical limit is therefore
$\Omega_c \simeq 17$--$18$\,rad\,s$^{-1}$, about a third lower than the
$\simeq 26$\,rad\,s$^{-1}$ reached in the field-free case. The
simultaneous, progressive failure of two independent diagnostics is the
signature of a genuine loss of virial balance rather than of the
solver jumping to a spurious branch. The $+47\%$ of
Table~\ref{tab:masses} should therefore be read as a static ceiling for
the toroidal branch on its own, not as a floor to which rotational
support can simply be added.

%-------------------------------------------------------------------

\section{Limitations}
\label{sec:limitations}

Three limitations bound what the results above can be taken to mean,
and none of them is incidental to the method.

\paragraph{No stability test.} These are two-dimensional axisymmetric
equilibria, certified to satisfy hydrostatic and virial balance. A
purely toroidal field is precisely the configuration subject to the
Tayler $m=1$ instability \citep{tayler1973,markey1973}, which is
non-axisymmetric and therefore invisible to this formulation by
construction. Nothing here tests whether the configurations of
Table~\ref{tab:masses} survive a non-axisymmetric perturbation, and the
expectation from stability theory is that, in this pure form, they do
not. A companion investigation addresses this by relaxing a random
seed field dynamically in three-dimensional MHD and measuring which
mixed poloidal-toroidal configuration survives, rather than imposing
one.

\paragraph{No exterior field.} By Eq.~\ref{eq:fieldlaw} the field
vanishes with the density, so it is confined to the stellar interior
and the exterior solution is field-free. There is no surface dipole,
and hence nothing for magnetic braking to act on -- which makes this
branch, on its own, unable to represent the rotational evolution of a
magnetized white dwarf, however well it represents its equilibrium
structure.

\paragraph{An unresolved corner of parameter space.} At
$K = 5\times10^{-3}$ mesh refinement does not converge cleanly: VE
falls from $8.2\times10^{-3}$ to $9.6\times10^{-4}$ and then rises
again to $3.7\times10^{-3}$ at the finest mesh tested, a
non-monotonicity that survives the Green's-function fix of
Sect.~\ref{sec:scf} and that we have not been able to attribute either
to a genuine sensitivity near a marginal equilibrium or to a further
numerical effect. We flag it as open rather than resolve it in either
direction. It does not affect $K = 3\times10^{-3}$, the value carrying
the results of Sect.~\ref{sec:results}, which converges monotonically
under mesh refinement ($1.5\times10^{-2} \rightarrow 1.7\times10^{-3}
\rightarrow 2.7\times10^{-4}$) and is unchanged to four significant
figures under domain enlargement.

%-------------------------------------------------------------------

\section{Conclusions}

TBD.

\begin{acknowledgements}
TBD.
\end{acknowledgements}

%-------------------------------------------------------------------

\bibliographystyle{aa}
\bibliography{paper}

\end{document}
